\begin{papercard}{hsieh2005}
\cf{Stages} \stagetag{C} (also \stagetag{L}). The classic typology of qualitative content analysis: \emph{conventional} (inductive), \emph{directed} (theory-driven) and \emph{summative} (keyword counting plus context). \emph{Note:} PDF page 1 is the SAGE cover sheet, so journal page 1277 is PDF page 2. The examples are hypothetical end-of-life studies.

\cf{Definition and scope} Qualitative content analysis is ``a research method for the subjective interpretation of the content of text data through the systematic classification process of coding and identifying themes or patterns'' \pc{hsieh2005}{3}{qualitative content analysis is defined as a research method for the subjective interpretation}. It goes beyond counting words to classify large amounts of text into an efficient number of categories with similar meaning, explicit or inferred \pc{hsieh2005}{3}{Qualitative content anal- ysis goes beyond merely counting words}. The three approaches differ in \emph{coding scheme, origin of codes and threats to trustworthiness} \pc{hsieh2005}{2}{The major differences among the approaches are coding schemes, origins of codes, and threats to trustworthiness}.

\cf{Conventional content analysis} Aim: describe a phenomenon where theory or literature is limited. Categories flow from the data (inductive category development). Steps: read all data repeatedly for immersion; read word by word and highlight exact words capturing key concepts; write first impressions; let code labels emerge, often in-vivo; sort codes into categories (ideally 10--15 clusters); build a tree diagram; define each category, subcategory and code; pick exemplars \pc{hsieh2005}{4}{Data analysis starts with reading all data repeatedly to achieve immersion}. Theory enters in the discussion. Table 1 example: open-code three to four transcripts, fix preliminary codes, recode, add codes when needed, then merge or split. Strength: direct information without imposed categories. Challenge: incomplete contextual understanding, i.e.\ missing key categories (a \emph{credibility} threat), addressed by peer debriefing, prolonged engagement, triangulation, negative case analysis and member checks \pc{hsieh2005}{5}{Credibility can be established through activities such as peer debriefing}. It is limited to concept development or model building \pc{hsieh2005}{6}{At most, the result of a conven- tional content analysis is concept development or model building}.

\cf{Directed content analysis} Aim: validate or extend a theoretical framework (``deductive category application''). Steps: identify key concepts or variables from theory or prior research as initial codes; define operational definitions from the theory; code with predetermined categories. Two strategies: (1) highlight all relevant text first, then code with the predetermined codes; (2) start coding immediately with the predetermined codes. Uncodable text gets a new code or is analysed later as a possible new (sub)category \pc{hsieh2005}{6}{Coding can begin with one of two strategies, depending on the research question}, \pc{hsieh2005}{7}{Data that cannot be coded are iden- tified and analyzed later}. Findings: supporting and non-supporting evidence, rank-order frequency of codes, share of supporting vs.\ non-supporting codes. Limits: informed but strong bias toward confirming the theory, leading probes cue participants, and theory can blind the researcher to context. Remedy: an \emph{audit trail} and an auditor reviewing operational definitions before the study \pc{hsieh2005}{8}{To achieve neutral or unbiased results, an audit trail and audit process can be used}.

\cf{Summative content analysis} Starts with identifying and quantifying words or content to understand how they are used; the counting explores usage and does not infer meaning. If it stops there it is quantitative (\emph{manifest}); it becomes qualitative when it interprets latent content and contexts \pc{hsieh2005}{8}{A summative approach to qualitative content analysis goes beyond mere word counts to include latent}. Procedure: computer or hand searches for terms, counts by source or speaker, then identify alternative terms (euphemisms) and compare usage across variables. Used for analysing manuscript types or textbook content. Strengths: unobtrusive, shows actual use. Limits: inattention to broader meanings. Trustworthiness: credibility, shown by expert validation of the term list or member checks \pc{hsieh2005}{10}{For Researcher Z’s study, validation by content experts on what terms are used to replace}.

\cf{Common process and the coding scheme} Seven classic steps (Kaid): research questions, sample, categories, coding process and coder training, implementation, trustworthiness, analysis of the coding \pc{hsieh2005}{10}{All approaches to qualitative content analysis require a similar analytical process of seven classic steps}. A coding scheme is a ``translation device'' with systematic rules, and its development is central to trustworthiness. Table 4 contrasts the starting point (observation, theory, keywords), the timing of defining codes (during analysis, before and during, keywords before and during) and their source (data, theory or research, researchers' interest or literature) \pc{hsieh2005}{11}{Key differences among conventional, directed, and summative approaches to content analysis center on how initial codes are developed}.

\cf{Conclusion} Choose the approach by matching research purpose to the state of the science. Name the approach \emph{before} analysis begins, because a stated procedure or coding scheme increases trustworthiness \pc{hsieh2005}{11}{It is important for health researchers to delineate the specific approach to con- tent analysis they are going to use}.

\cf{Use for this thesis} A vocabulary for the codebook design. The thesis codebook is most naturally \emph{directed}: initial codes come from prior work (Malki, Hutton, Habib, dark-pattern taxonomies, GDPR transparency criteria), with an explicit rule for uncodable observations (new code vs.\ ``other'', reviewed in the pilot). The \emph{summative} approach fits counting terms in policies (e.g.\ the words ``anonymous'' or ``may'') followed by contextual interpretation. State which approach each analysis step uses and the bias risk that comes with it (directed designs favour confirmation), and use an audit trail (decision log, versioned codebook) as the mitigation Hsieh \& Shannon propose. The paper offers no reliability procedure, so pair it with MacQueen and O'Connor \& Joffe.
\end{papercard}
